\part{Introduction}

\chapter{Project Overview}
\label{chap:introduction}

Cardiovascular disease remains the leading cause of mortality worldwide. The electrocardiogram (ECG) serves as the primary non-invasive diagnostic tool for arrhythmia detection, yet manual analysis of long recordings is error-prone and time-consuming. Deep learning models have demonstrated cardiologist-level performance in ECG classification tasks, offering the potential to augment clinical decision-making.

\section{Problem Statement}
Traditional ECG analysis relies on manual feature extraction and shallow machine learning models, which fail to capture complex morphological and temporal patterns in ECG signals. This project addresses the need for automated, accurate, and generalizable arrhythmia classification using state-of-the-art deep learning architectures.

\section{Project Objectives}
The main objectives of this mini-project are:
\begin{enumerate}
    \item Implement and compare 4 deep learning architectures for 5-class ECG arrhythmia classification
    \item Develop a unified, reproducible training pipeline with logging and checkpointing
    \item Deploy the best-performing model in a user-facing web application
    \item Align with course requirements: CNNs, RNNs, Transformers, Transfer Learning, Dimensionality Reduction
\end{enumerate}

\section{Scope and Constraints}
The project adheres to the following constraints:
\begin{itemize}
    \item \textbf{Allowed Architectures:} 1D CNNs, RNNs, Transformers, Autoencoders (no text as primary modality)
    \item \textbf{Dataset:} MIT-BIH Arrhythmia Database (publicly available, standard benchmark)
    \item \textbf{Deliverables:} Web application, detailed report, GitHub repository
    \item \textbf{Course Alignment:} Complies with Deep Learning module content.
\end{itemize}

\section{Report Structure}
This report follows standard academic structure: introduction, related work, dataset details, methodology, implementation, results, deployment, discussion, conclusion, and references.
